% =======================================
% MEETING 6
% ---------------------------------------
% DATE   - SEPTEMBER 4, 2026
% FORMAT - ONLINE
% =======================================
% TOPIC  - Vector transformations, Curves and Tangent Vectors 
% =======================================

\subsection{TBA}
\subsubsection{Structures in Manifolds}

Last meeting, we talked about structures in manifolds.
These are the mathematical objects that will be important
in general relativity. We discussed scalars, which was 
fairly easy, and we also started looking at vectors. 

The key thing here is that a vector on a manifold is really
a map, it's a derivative operator, or a derivation, at least
this is how we tried generalizing it. 

So by that, we mean it's a map that can act on functions. So
if you have a vector $\vb{v}_p$ on a point $p\in\mcal M$, which belongs to $T_p \mcal M$,
\[
	\vb{v}_p \in \TpM 
\]
that means it can act on functions, and these actions on functions, will
be a real number.
\[
	\vb{v}_p: f(\mcal M) \to \R
\]
So it's kind of like a directional derivative. The 
intuition is that if you provide a function to $\vb{v}_p$,
it will return a rate of change on that function along some kind 
of direction. 

We haven't really completely filled up the details of that 
story. As we discussed before, we have to recover the arrow again. We 
ask how do we explicitly make a vector look like a directional 
dirivative in $\R^n$, so that's in part what we are going to
achieve today. 

The other thing we discussed last time is that $\TpM$, otherwise called 
a tangent space (we will justify why that name is appropriate), is basically
the set of all vectors, or directional derivatives, at $p$. It's a vector space or a 
linear space, but for it to become a vector space we had to define what the 
operation of plus, $+$, means, what scalar multiplication means, and this is 
what we introduced also in the last lecture. 

Now, being a vector space, one of the most important things one can do when
one has a vector space is establish a basis on that vector space. We know already
from our many experiences with vectors spaces in physics. One nice basis is the so-called
coordinate basis on $\TpM$. The coordinate basis is one that we get from having a 
chart on $p$.

So if we have a manifold and a point $p$, remember that it's going to be an open
set that contains $p$ and then there's going to be a chart $\psi$ that maps $p$ to a coordinate system.
\begin{figure}[h]
    \centering
    \includegraphics[width=3in, height=2in]{psimapping.pdf}
    \caption{}
    \label{fig:psimapping}
\end{figure}


This, we often just write as $\{x^\alpha\}$, instead of saying $\psi$, which is kind of fairly 
abstract, we just say $\qty{x^\alpha}$, we express it in terms of the coordinate functions.
$\qty{x^\alpha}$ are the coordinate functions that manifest on psi.
\[
	\psi \iff \qty{x^\alpha}
\]
and once we have that we have our basis, which is expressed as a directional derivative
in terms of the coordinate functions.
\[
	\qty{\pdv{}{x^\alpha}}
\]
Often times, we will write this simply as 
\[
	\qty{\p_\alpha}
\]
When it's confusing, for example you have two different baes, we will write this differently.
For example you have a different basis from a different chart $\qty{y^\alpha}$, we instead
explicitly define each basis vectors
\[
	\psi_x \implies\qty{\pdv{}{x^\alpha}},\qquad \psi_y \implies \qty{\pdv{}{y^\alpha}}
\]
An example would be different representatiosn of a vector in cartesiaon coordinates in $\R^3$
or spherical coordinates in $\R^3$. As a shorthand, we can also write
\[
	\qty{\pdv{}{x^\alpha}}=\qty{\partial_\alpha},\qquad \qty{\pdv{}{y^\alpha}}=\qty{\partial'_\alpha}
\]
If there is confusion, like when there's a possibility that we're working on different charts, with 
two different sets of coordinate vectors, we can use this convention of having a primed operator.

We also discussed last time that since it's a basis, I can now express any vector $\vb{v}_p$ in terms
of any coodinate basis, so we can write, in terms of the $x$ chart,
\[
	\vb{p} = 
	v_p^\alpha \partial_\alpha = 
	v_p^\alpha \pdv{}{x^\alpha}  
\]
Here , we said that $v_p\alpha$ is really nothing but the actions of $\vb{v}_p$ on the 
coordinate basis function $x^\alpha$. This is what we established last time. 
\[
	v_p^\alpha = \vb{v}_p(x^\alpha)  
\]
\begin{question}{So is $\partial^\alpha$ equivalent to the elementary basis vectors $\vb{e}^i$?}
	No, not all the time. They are kind of similar but they are not quite. The elementary basis vectors
	in traditional physics, are equivalent to it in cartesian $\R^3$, simply
	a different way of writing
	\[
		\ihat, \jhat, \khat \equiv \pdv{}{x},\pdv{}{y},\pdv{}{z} 
	\]
	However, in other systems of coordinates, not necessary, because when we say
	something like $\vb{e}^\alpha$, then we are already saying something about the magnitude
	of these vectors, since these are unit and orthogonal vectors, whereas $\partial^\alpha$ 
	are not necessarily unit nor orthogonal because we do not have a notion of length here yet, 
	and we do not have a notion of orthogonality as well. Those statements precedes those notions,
	We can have coordinate basis vectors even without the notion of lengths or orthogonalities or 
	inner products. In cartesian, these will turn out indeed to be equivalent once we introduce 
	the so called Euclidean metric. The Euclidean metric is what's going to provide us a sense of 
	length and orthogonality which is an inner product. Prior to that, without the metric, we already
	have coordinate basis vectors, but not the usuak sense of lengths or orthoginality, which $\ihat,\jhat$, and $\khat$
	already containts.  

	\hspace{1em}For example, the spherical basis vectors, they point in the same direction as $\partial_\alpha$, but they
	are not equivalent to the coordinate basis vectors you get from the manifold.
	\[
		\vu*{r}, \vu*{\theta}, \vu*{\phi}\not\equiv \pdv{}{r},\pdv{}{\theta}, \pdv{}{\phi}  
	\]
	These are the still coordinate basis vectors, it's very easy you just take the coordinate and put it on the denominator,
	but they have (you will see this when we introduce the meaning of length) different
	magnitudes. Later on, when we introduce the notion of normalization and define a unit vector, then we recover the same vectors. 
\end{question}

Now, we proceed to the question that we left hanging in the last lecture. 

Suppose $p\in\mcal M$ belongs to two different charts, $\psi_1$ and $\psi_2$, and as such
we have two different sets of coordinate functions that apply and describe $p$.

If we that is the case, then we also have different sets of coordinate basis vectors on $\TpM$. It's 
the same $\TpM$, but for each chart, it can be described by two different sets of basis vectors, $\qty{x^\alpha}$ and $\qty{y^\alpha}$. 
So the question becomes, how do the different sets of components of $p$ relate from one chart to another? Note that
the components $v^\alpha$ is not the same as the vector $\vb{v}$. The components can change
but the vector remains the same. 
\begin{figure}[h]
    \centering
    \includegraphics[width=4in,height=3in]{twocharts.pdf}
    \caption{}
    \label{fig:}
\end{figure}

Notice that $x^\alpha$ is the $\alpha$th coordinate basis on the chart $\psi_1$, since $\psi_1:p\mapsto x^\alpha$ 
\[
	x^\alpha = \psi_1^\alpha(p)
\]
Similarly for $\psi_2$, we have
\[
	y^\alpha = \psi_2^\alpha(p)
\]
We asked last lecture, how do we express a coordinate transformation from $y$ to $x$ in terms of $\psi_1$ and $\psi_2$. We find that
\[
	y^\alpha (x^\beta) = (\psi_2 \circ \psi_1\inv)^\alpha(x^\beta)
\]
Now we want to understand how do we define a vector present in both charts, and how does its component relate to each other. 

To specify a vector, the way we have defined them last lecture, is that it's a map from functions to real numbers. We first need to define the map.
Every vector is a derivative operator. If you tell me how it operates on the function, then you will have specified the vector. 

Suppose a vector on the $y$ basis at point $p$ acting on a function. Recall how we defined this last lecture
\[
	\qty( \pdv{}{y^\alpha})_p f = \eval{\pdv{f\circ\psi_2\inv}{y^\alpha}}_{\vb{y}_p}
	\label{eq:y-vecotr}
\]
Similarly, the $x$ coordinate vectors have exactly the same form.
\[
	\qty( \pdv{}{x^\alpha})_p f = \eval{\pdv{f\circ\psi_1\inv}{x^\alpha}}_{\vb{x}_p}
\]
Notice that the definitions are exactly the same, the only difference is that we the corresponding chart for each coordinate system.
So, how do these things relate to each other? That's what we're trying to address.

We can insert, in essence, an identity map $\psi\inv\circ\psi$
\[
	\qty( \pdv{}{x^\alpha})_p f 
	= \eval{\pdv{f\circ\psi_2\inv \circ \psi_2 \circ\psi_1\inv}{x^\alpha}}_{\vb{x}_p}
\]
And then, we can just expand this out with standard multivariable calculus techniques, in particular 
we use the chain rule.
\begin{question}{Is $\psi\circ\psi\inv$ also an identity?}
    Yes, it is also an identity map. However, we had to do it in this particular order for this derivation
	because we will not get the result we desire due to the non-commutative nature of 
	these operators. 
\end{question}

Proceeding with the operation, we find that
\[
	\pqty{ \pdv{}{x^\alpha}}_p f 
	=
	\eval{\pdv{(f\circ\psi_2\inv)}{y^\beta}}_{\vb{y}_p} \eval{\pdv{(\psi_2\circ\psi_1\inv)^\beta}{x^\alpha}}_{\vb{x}_p} 
	\label{eq:id-operator}
\]
Notice the change to $\vb{y}_p$. Since $f\circ \psi_2\inv$, as defined, is simply $f(\vb{y}_p)$, we evaluate from there.
Now, wat is $\psi_2\circ\psi_1\inv$? From the last lectures, where you can verify, this is just $\vb{y}(\vb{x})$, i.e., it
can be represented as $y^\beta$.

Also notice that we used another dummy index $\beta$. Since we used the chain rule, which sums over intermediate coordinates,
we need to represent the intermediate coordinate index separately from the fixed 
$\alpha$, which labels the $x$-coordinate we are differentiating with respect to.

Now we can represent (\ref{eq:id-operator}) as 
\[
	\qty( \pdv{}{x^\alpha})_p f 
	=
	\eval{\pdv{(f\circ\psi_2\inv)}{y^\beta}}_{\vb{y}_p} 
	\eval{\pdv{y^\beta}{x^\alpha}}_{\vb{x}_p} 
\]
But now look at it, we see that the first term is just (\ref{eq:yvector}). This is nothing but
\[
	\pqty{ \pdv{}{x^\alpha}}_p f 
	=
	\pqty{\pdv{}{y^\beta}}f \cdot 
	\eval{\pdv{y^\beta}{x^\alpha}}_{\vb{x}_p} 
\]
We can rearrange this, because this is just multiplication, and $\displaystyle\eval{\pdv{y^\beta}{x^\alpha}}_{\vb{x}_p}$ are just numbers, then
\[
	\pqty{ \pdv{}{x^\alpha}}_p f 
	=
	\eval{\pdv{y^\beta}{x^\alpha}}_{\vb{x}_p} 
	\pqty{\pdv{}{y^\beta}}f  
\]
Notice that this is for an arbitrary function $f$, and as such, in an operational sense, we can write
\[
	\pqty{ \pdv{}{x^\alpha}}_p 
	=
	\eval{\pdv{y^\beta}{x^\alpha}}_{\vb{x}_p} 
	\pqty{{\pdv{}{y^\beta}}}_p
	\label{eq:chain-rule-result}
\]

One of the things that you will obviously see here is that this is nothing but the expression 
of the chain rule in multivariable calculus. It's as if we just wrote down the chain rule. It's very 
easy to remember. However, note that this has a different function, we are sort of giving a new interpretation
to the chain rule. Becayse what this statement actually telling you? 
What we're doing here is expressing the left hand side, which is a vector, as a superposition or a sum over other 
basis vectors. So we are expressing a vector on the left hand side as a linear combination of vecotrs on the right hand side. 

So something like the chain rule, which is taken for granted in multivariable calculus because it is used so 
often, now gets a different interpretation in terms of vectors. The values on (\ref{eq:chain-rule-result})  are 
just numbers, so it's as if we are doing
\[
	\ket{\psi} = \psi^i \ket{e_i}
\]
or a matrix multiplcation 
\[
	\vb{x} = \vb{M}\vb{y}
\]

We interpret (\ref{eq:chain-rule-result}) as the following: The $\alpha$th basis vector in the $x$-coordinate system
is a linear combination of the basis vectors in the $y$-coordinate system. Similarly, you can express
\[
	\eval{\pdv{}{y^\alpha}}_p =  
	\eval{\pqty{
		\pdv{x^\beta}{y^\alpha} 
	}}_{\vb{y}_p} 
	\pqty{{\pdv{}{x^\beta}}}_p
\]
So it means, that we can represent one basis in terms of the other. So how then do 
we relate the different components comming from these basis vectors? 

What we did earlier is that given two sets of basis vectors, how do we relate one in terms 
of the other?

By the way, the expression $\pdv{y^\beta}{x^\alpha}$ are Jacobians, this is often called the tangent map,
we will explain later on. This is actually the partial derivatives of the transition functions that are defined 
to be smooth in smooth manifolds. 

If I have a vector$ \vb{v}_p \in \TpM$, and I have two different charts, then I get 
to represent this vector using two different sets of basis vectors. I can say that
\[
	\vb{v}_p = v_p^\mu \pdv{}{x^\mu}
	\label{eq:vpx}
\]
or we can write the same vector as 
\[
	\vb{v}_p = {v'_p}^\nu \pdv{}{y^\nu}
	\label{eq:vpy}
\]
The important thing here is that we are talking about the same vector. Just like we can represent the vectors
in different ways in different coordinate systems
\begin{align*}
	\vb{v} &= v_x \ihat + v_y \jhat + v_z \khat\\
	\vb{v} &= v_r \vu*{r} + v_\theta \vu*{\theta} + v_\phi \vu*{\phi}
\end{align*}

Since they'are the same vector, then we can equate (\ref{eq:vpx}) and (\ref{eq:vpy})
\[
	{v'_p}^\nu \pdv{}{y^\nu}
	=	
	{v_p}^\mu \pdv{}{x^\mu}
\]
But, we can express $\pdv{}{y^\nu}$ in terms of the coordinate basis vectors in $x$, as we have shown
\[
	{v'_p}^\nu
	\pqty{
		\pdv{x^\mu}{y^\nu}
		\pdv{}{x^\mu} 
	}
	=	
	{v_p}^\mu \pdv{}{x^\mu}
\]
Though remember, the mnemonic is that of the chain rule, it's functionally different, it's a linear combination of the basis bectors.
Proceeding, we can move the LHS to the RHS, and we find
\[
	\pqty{
		v'_p \pdv{x^\mu}{y^\nu} - v_p^\mu
	}
	\pdv{}{x^\mu}
	=
	\vb{0}
\]
Now we invoke the fact that we proved last lecture that $\pdv{}{x^\mu}$ is 
a linear basis, so they are linearly independent. So if the sum of the 
coefficients and the vectors are the zero vector, then the components inside the 
parenthesis are equal to zero. Hence
\[
	{v'}_p^\nu \pdv{x^\mu}{y^\nu} - v_p^\mu
	=
	0
\]
As such we can write
\[
	{v_p}^\mu = 
	{v'}_p^\nu \pdv{x^\mu}{y^\nu}  
	= 
	\pdv{x^\mu}{y^\nu}  {v'_p}^\nu
	\label{eq:vpp-to-vp} 
\]
Similarly, we can express (\ref{eq:vpp-to-vp}), by inverting the 
argument and taking the reciprocal of the derivative as
\[
	{v'_p}^\nu = 
	\pdv{y^\nu}{x^\mu} v_p^\mu 
	\label{eq:vp-to-vpp}
\]

So these are how the comonents of the vector $\vb{v}_p$ relate to one another under
a change of basis; this is a component transform.

Some of you will recognize that this is the Arfken\cite{Arfken1970} standard definition of a vector. In fact, it is a contravariant vector.
This is the old fashioned definition of a contravarient vector. This is what's being taught by Arfken and a lot 
of other elementary or introductory physics books. We say old fashioned because this is how it is traditionally done 
before,  but for us it is insufficient. For most physicists and most physics work, it is probably sufficient, but it doesn't 
really answer what a vector is, in a rigorous sense.

If you ask ``what is a vector", a physicist would say, ``Well, if its components transform this way, it's a vector" or ``It's an element of a vector space'', without
really saying or actually answering what a vector is. So now, we know what a vector is, it's a derivative
operator, and as a consequence of that, its components transform in a manner described by (\ref{eq:vp-to-vpp}).

So this is no longer the definition of a contravariant vector, it is the transformation rule arising from the fact that
vectors are derivative operators. For classical physics, it is generally sufficient to say how it transforms, but the problem is 
in GR and in fact for most of physics we are emphasizing the geometric object itself, defining
what is the object ranther than how its components transform under different coordinate systems. 

We are trying to divorce ourselves from coordinates as much as possible. We are trying to work with the geometric objects themselves.
Now we say it is a contravariant vector because there is a different kind of vector which is a covariant vector.

If you have a covariant vector $\omega_p$, the transformation rules are expressed as 
\begin{equation}
	\label{eq:covariantvectordef}
	(\omega_p)'_\nu = 
	\pdv{x^\mu}{y^\nu} \pqty{\omega_p}_\mu 
\end{equation}
This is the definition of a covariany vector. But later on, in a couple lectures from now,
we will see that we will no longer call these vectors. These are actually going to be one forms.
Something which transforms like (\ref{eq:covariantvectordef}) is called a one form, or alternatively, a 
dual vector. 

So that we aren't so confused, we have encountered these in physics before, a contravariant vector is equivalent to a \textit{ket},
\[
	\ket{\psi}
\]
and a covariant vector is equivalent to a \textit{bra}
\[
	\bra{\phi}
\]
We do not say that bra's and ket's are the same. Of course, they are isomorphic, but they
are quite different. They serve different purposes. 
Covariant and contravariant vectors are usually defined in terms of how they
transform in physics, but we will define these other vectors later on, when we discuss one forms.

Now it's importnat to emphasize that $\TpM$ is a vector space. You don't always need coordinate basis vectors for vector spaces.
In fact, later on we will have find that it is quite practical and useful to adopt non-coordinate basis vectors. 

So the vector basis need not be coordinate basis vectors, in other words, they don't need to come from a coordinate system or a chart 
because it's a space, you can put in a plethora, an infinite number of basis vectors on it. 

For example, you have a set of basis vectors, you will need $n$ of them if you have an $n$-dimensional manifld, and
they will have some components with respect to some basis, and these basis need not come from 
a coordinate system, it can be a lienar combination of coordinate basis vectors.
\[
	e_{(\beta)} = e_{(\beta)}^\alpha \pdv{}{x^\alpha} 
\]
The only constraint is 
that it must have a nonzero determinant. You don't get a basis unelss the matrix formed by the basis vectors if 
the matrix isn't a singular matrix. It has to be singular, otherwise you don't get a one-to-one map from the manifold to $\R^n$.

Later on when we go to Lie derivatives, we will investigate under what conditions $e_{(\beta)}$ will actially have a coordinate representation.
The short answer is that the Lie derivative should be zero
\[
	\mathscr{L}_{(\beta)} e_{(\alpha)} = 0
\]
If this is true for any $\alpha$ and $\beta$ for this new set of basis vectors that we've introduced, then 
there is a set of coordinates for which these are coordinate basis vectors. I can find a coordinate system to write
for example
\[
	e_{\beta} = \pdv{}{z^\beta} 
\]
If the Lie derivative condition is satisfied. But we'll get to that much later on. 

So now, we want to be able to recover the notion of an arrow in $\R^n$. So far, we've had a lot of properties 
sort of found and recovered in what we understand to be traditional vectors in physics. But to really get 
a sense of an arrow, we need to talk about curves and tangent vectors.

\subsubsection{Curves and Tangent Vectors}

We feel like we have a notion or sense of what a curve is, since we encounter it in physics 
a lot. However, for our purposes, we need to precisely define what a curve is.
\begin{definition}{Curves}
	A curve $C$ in a manifold $\mcal M$ is a map from an open interval $I\subset\R$ 
	onto $\mcal M$
	\begin{equation}
		\label{eq:curvedef}
		C:I \to\mcal{M}
	\end{equation}
\end{definition}
\begin{figure}[h]
    \centering
    \includegraphics[width=4in,height=2in]{curve_in_a_manifold.pdf}
    \caption{A curve in a manifold}
    \label{fig:curve-in-a-manifold}
\end{figure}

$C(I)$, is a function, which is different from how we usually imagine a curve, which is 
some line with a curvature on $R^n$. $C(I)$ is not actually the curve itself, but 
the image of the curve. So what we tend to think of the curve is really the image 
of the curve $C$, because $C$ itself is the map, $C$ is the function. It's the instruction manual that tells
us how to relate a real number to a point in $\mcal M$. $C(I)$ is a set of points in $\mcal M$. 
And $I\subset \R$ is called a parameterization. $I$ need not be a propersubset, it can be the entire real line. It's how we parameterize a curve. Say for example, you want to think of a parameter
$I$ is time, $t$, which is a very typical parameter in physics, the parameter need not be uniform. People can have different parameterizations,
different clocks running at different rates for instance. 

Imagine a different curve, say $C':I'\to \mcal{M}$, with different parameters $I'$, such that
\[
	C(I)= C'(I')
\]
then we call these curves have the same image, but different parameterizations. So we often say
that $C$ and $C'$ are reparameterizations of each other. It has the same image, the same set of points,
but the parameter $I$ that maps it is different, i.e., $C'(t')\neq C(t)$.

If they are reparameterizations of each other, the it is not so stramge that there will exist an onto map,
a subjective map $\alpha: I \to I'$, such that
\[
	\alpha(t)=t'
\]
and 
\[
	C = C'\circ \alpha
\]
An important thing to emphasize here is that a curve on its own does not need coordinates.
It's a map from the subset of the real line to the manifold. However, in physics, it's very common
to express curves in terms of coordinates, which we can do if we are on manifolds, mapping the points
$C(I)$ to $\R^n$. Pictorially, we 
can represent this as 
\begin{figure}[h]
    \centering
    \includegraphics[width=4in,height=2in]{curves-chart.pdf}
    \caption{$\psi:C(I)\to \R^3$}
    \label{fig:cuves=chart}
\end{figure}
As such, we can recover the normal parametrization of a real number as a coordinate in $\R^n$.
\[
	\psi\circ C: I \to \R^n
\]
in other words
\[
	\psi\circ C: t \mapsto x^\alpha(t)
\]
So for example, $[x(t),y(t),z(t)]$, which is how we normally describe curves in physics. 
You can provide a coordinate representation of curves but the curve itself actually did not rquire 
the coordinates. So what we do is coordinatize a curve. 

There are goign to be these special curves which we need to define. Given some chart, which we represent with $\qty{\alpha}$,
consider the following subsets of open sets in $\mcal M$ for which the coordinates are valid, with chart $\psi$, $\qty{\mcal{O},\psi}$. Consider the following
subset of all points in $\mcal{O}$ such that 
\[
	\qty{
		p\in\mcal{O} : 
		x^2(p) = \text{constant ,}
		x^3(p) = \text{constant ,}
		\ldots
		x^n(p) = \text{constant }
	}
\]
What do you imagine these sets of points to be? The image of this map is a straight line at $x^1(p)$, because we are keeping all the $x^\alpha$'s constant, except
for $x^1$. However, the set that we're looking at, since it's a set of all $p$'s in $\mcal O$, is the pre-image of the 
points that will get mapped into a line. So this set is asking you to imagine what points in $\mcal O$ will get mapped into a 
straight line by $\psi$

\begin{figure}[h]
    \centering
	\includegraphics[width=4in,height=2in]{specialcurves.pdf}
    \caption{}
    \label{fig:specialcurve}
\end{figure}

So with these set, consider the inverse map $\psi\inv$, it defines a set of points that will get mapped to a straight line. Since we konw it is contunous,
then that necessarily means that the pre-image is a continuum of arbitrary points, i.e., the preimage is a curve on the manifold, because it's mapping an 
open interval of $\R$ onto a set of points in $\ncal M$. This is a special type of curve, it's called the $x^1$ coordinate line or $x^1$ coordinate curve. 
Notice here that the parameter is $x^1$. This is an important notion.

The $x^1$ coordinate curve is the image of a curve on the manifold that maps the parameter we call $x^1$ coordinate.

Related to curves are things called tangent vectors. 
\subsubsection{Tangent vectors}

You have an intuitive sense of what is a tangent vector, it's the arrow that points in the direction
instantantously of the curve. IN a sense, a tangent vector is our traditional notion of what a tangent vector is.
But because it's a vector, we now know that it's supposed to be a directional derivative. 
We have to apppropriately define then, a tangent vector, asa derivative operator. 

\begin{figure}[h]
    \centering
	\includegraphics[width=4in,height=2in]{tangentvectors.pdf}
    \caption{}
    \label{fig:tangentvectors}
\end{figure}

It's always associated with curve, i.e., if there's a curve, there must be a tanget vector. We define it then as usch.

\begin{definition}{Tangent Vector}
	If $C(t)$ is a curve, then the tangent vector $\tau$ at $C(t_0)$ is the derivative operator
	such that $\tau_{C(t)}$ acting on an arbitrary function $f$, is defined to be
	\begin{equation}
		\label{eq:tangent_vector}
		\tau_{C(t_0)} (f) = 
		\eval{\dv{(f\circ C)}{t}}_{t_0}
	\end{equation}
\end{definition}

What is the map $f\circ C$. Remember that $f$ is a map from the manifold to real numbers, and $C$ is a map 
from real numbers to the manifold. SO $f\circ C$ is a map from real numbers to real numbers. 
\[
	f\circ C : \R \to \R
\]
So it's 
simply a normal function, a perfectly sensible function, which is a derivative
with respect to a parameter. 

Notice that the definition of a tangent vector requires a curve. If you have the same curve $C$,
but evaluating another parameter $t_1$, then it makes a different tangent vector acting on
the same function, but evaluated at a different point
\[
	\tau_{C(t_1)} (f) = 
	\eval{\dv{(f\circ C)}{t}}_{t_1}
\]
This is a different tangent vector. We can show that this is in fact a vector as we previously defined it, i.e., 
it must satisfy linearity, and the Leibniz propety, but by the nature of the function, since it's a derivative operator, it obviously 
satisfies both. This definition will satisfy the properties of being a vector, or being a derivation. We will recover the notion
of an arrow from this, this is where the traditional arrow definition of a vector comes from. 













